% !TeX root = ../main.tex
\xchapter{统一可逆生成框架}{Crossreferences}
\xsection{引言}{Introduction}
生成对抗网络~\cite{goodfellow2020generative}通过对抗训练生成图像，但该范式在训练过程中往往引入不稳定性，且难以稳定扩展到更大规模。另一方面，受Simoncelli E. P.等早期在噪声建模与统计物理相关研究的启发~\cite{simoncelli1996noise,hyvarinen2005estimation,sohl2015deep}，生成模型（Diffusion Models）在DDPM~\cite{ho2020denoising}与基于分布分数的SBDM框架~\cite{song2021maximum,song2020score}中被系统化提出并迅速发展。
从经典随机过程视角，生成模型常可概括为流模型、扩散模型与桥模型三类。尽管三者在随机过程的具体构造与训练形式上存在差异，但流模型与扩散模型在本质上都利用神经网络刻画一个时间连续的随机动力系统${x_t}$：其训练目标是学习由简单先验分布$p_0(\epsilon)$向数据分布$p_1(x)$演化的概率流，并通过求解该随机过程的逆向随机微分方程（reverse-time SDE）或等价的概率流常微分方程（probability flow ODE），实现从噪声到数据样本的可采样映射。相比之下，桥模型可视为在给定初末分布约束下的条件化流模型，即在保持端点分布一致性的前提下，对上述概率流过程进一步施加路径层面的约束，从而实现受控的分布迁移。
在此基础上，研究社区通常从提升神经网络的建模能力~\cite{ho2020denoising,nichol2021improved,rombach2022high,peebles2023scalable,bao2023all}与改进生成路径~\cite{lipman2022flow,peluchetti2023non,tong2023improving,shi2023diffusion,zhou2024denoising}两方面入手，以获得更稳定且更高保真的生成结果。


在提高神经网络的建模能力方面，
Jain等人~\cite{ho2020denoising}提出的 DDPM 工作提出将 UNet 作为时间条件去噪器的范式，即以对称的编码器与解码器结构在多尺度上提取特征，并通过长跳跃连接对齐浅层纹理信息与深层语义信息，同时将时间步嵌入注入到残差块中以建模不同噪声尺度下的去噪映射。该设计使模型能够在所有 $t$ 上对扰动分布 $q(x_t)$ 的去噪方向进行稳定学习，从而以逐步反演的方式逼近数据分布。随后，Nichol 与 Dhariwal~\cite{nichol2021improved}系统讨论了注意力层设计、噪声方差参数化与训练目标等因素对对数似然与生成质量的影响，Rombach 等人~\cite{rombach2022high}在潜空间扩散模型中将去噪过程转移到自编码器潜表示上，并在 UNet 中引入交叉注意力层以实现对文本与边界框等更一般条件的选择性融合，从而在降低计算成本的同时保持高分辨率细节生成能力，将扩散模型的生成质量推进到可与当时 GAN 竞争的水平。Peebles 与 Xie 提出的 DiT~\cite{peebles2023scalable}以 Transformer 替代传统 UNet 主干，在隐空间上执行去噪建模，并通过归一化调制等方式将时间步与条件信息注入到 Transformer 块中，展示了随模型深度、宽度与 token 数增长而持续收益的可扩展性，为扩散模型借鉴大模型训练范式与大规模数据训练提供了更直接的结构基础。进一步地，Bao 等人在 U-ViT~\cite{bao2023all}中提出 U 形模型主干并指出对扩散模型而言长跳跃连接是关键结构要素，而卷积 UNet 网络中的下采样与上采样算子并非总是必要，通过在浅层与深层 Transformer 块之间建立跨层特征通路可以改善收敛与细节生成质量，从而为探索更一般的扩散主干设计打开空间。总体上来看，以 Transformer 为主干的扩散网络正逐步与基于交叉注意力、归一化调制与旁路控制分支的条件注入机制相结合，使扩散模型在可扩展性与可控生成两方面同步增强。

在改进生成路径方面，Lipman 等人~\cite{lipman2022flow}提出用于图像生成的流匹配方法，其核心是直接学习连续时间速度场或等价的概率流，从而在采样时沿确定性或随机动力学将简单先验逐步推送到数据分布。通过采用合适的扩散桥混合，可以构造连接端点分布的扩散过程并实现对目标数据分布的生成~\cite{peluchetti2023non}。在此基础上，Tong等人~\cite{tong2023improving}引入广义条件流匹配技术，将外部条件视为对流场的约束或对端点的指定，使得条件信息能够以输入拼接、特征调制或注意力融合等方式注入参数化的主干网络中，并在训练阶段直接对齐条件下的路径分布。进一步地，针对更一般的桥建模需求，Shi 等人~\cite{shi2023diffusion}通过迭代式的马尔可夫拟合求解薛定谔桥的近似迭代，并以得分网络作为核心参数化模块，从而在同一主干上同时刻画中间时刻的去噪动力学与端点约束的桥一致性。Zhou 等人~\cite{zhou2024denoising}则将扩散桥显式作为生成与变换的统一建模对象，使用与扩散模型相同的得分网络主干来学习桥过程的得分函数，并将端点条件以输入端点样本、特征级注入或多尺度通路的形式并入网络，使得从给定源端点到目标端点的采样可通过求解相应随机微分方程直接完成。这些工作使生成路径不再局限于从噪声到数据的单端点设定，而是能够在更一般的端点约束与条件变换场景中复用生成模型的主干结构与训练范式，从而为生成模型向具身智能与基础科学等更广泛领域的拓展提供了更直接的技术基础。

然而，从上述两个方向可以看到，现有方法总体上仍以某一特定任务为目标，尚缺少能够同时覆盖流模型、扩散模型与桥模型的统一可逆生成框架。具体而言，仍主要围绕计算机视觉中的图像生成、视频生成与图像补全等单一任务展开，对生成模型在更一般机器学习问题及具身智能等下游场景中的迁移性与泛化性关注不足，因而难以形成面向多任务的共享表征与可复用训练范式。与此同时，当方法需要显式调用 ODE/SDE 等偏微分方程相关的数值求解器以实现更灵活的路径建模与约束注入时，现有方案往往依赖对网络结构、条件注入通路或离散化策略进行专门设计，导致模型主干与生成路径之间高度耦合，缺乏清晰的模块化分解与可复用接口，进而难以系统性支撑大规模预训练与向下游任务的高效迁移。更进一步，尽管概率流与逆向动力学在理论上蕴含可逆结构，许多方法并未充分挖掘这一结构来统一支撑信息重建、推断与可控生成等能力，在实践中，求解过程也常因不可微实现、截断近似或算子封装等因素，难以在 ODE/SDE 求解器层面实现端到端的梯度反向传播，从而限制了生成过程与下游目标的联合优化能力，阻碍了模型整体训练与任务微调。

因此，本文旨在从统一的视角出发，建立分数匹配与流匹配在模型建模方式与训练目标上的内在对应关系，并进一步刻画扩散路径与数值求解器之间的耦合机理。具体而言，本文一方面将不同扩散建模范式归一到一致的优化目标表述中，使得在同一主干网络参数化下，可以自由切换分数匹配或流匹配式用于建模或是训练，从而实现不同扩散方法与神经网络建模之间的系统化配合。
另一方面，本文将扩散路径视为连续动力系统的实例化结果，将求解器视为可替换的推断与训练组件，使得采样与反向传播能够在统一接口下完成，并支持通过替换不同求解器来平衡生成质量、稳定性与计算开销。
为进一步支持可逆扩散与端到端训练，本文引入 Hutchinson 迹估计器以高效近似连续动力学中的迹相关项，从而在不显著增加计算与内存成本的前提下获得可微的训练信号，并使基于求解器的梯度更新计算更可行、更稳定。
整体上，该框架在数学层面兼容多类扩散生成模型及其变体，同时在工程实现上支持以可插拔求解器完成训练与推断流程的统一。该框架受到社区的广泛关注，已在航天材料生成~\cite{yang2024data,zhang2025funcgenfoil}、图像合成~\cite{Xue_2025_ICCV}、强化学习~\cite{zhang2024generative_rl}等领域上获得应用与验证。

\xsection{生成模型框架}{Generative Model}
本文将首先介绍扩散模型与流模型的数学刻画，并讨论如何在统一的连续时间动力学框架下理解二者的建模与训练目标。

\xsubsection{扩散模型}{Diffusion Model}
令 $t$ 表示时间参数，记生成轨迹在时刻 $t$ 的随机变量为 $x_t$，其边缘分布记为 $p_t$。其中 $x_0 \sim p_{\mathrm{data}}$ 表示原始数据样本，也即原图；$x_1 \sim p_1$ 表示扩散终止时的先验分布，通常取为标准高斯 $\mathcal{N}(0,I)$。

扩散模型的核心思想是通过学习反向动力学来逆转一个前向扩散过程。前向加噪在条件分布层面可写为
\begin{equation}\label{eq:forward_diffusion}
p(x_t \mid x_0) \sim \mathcal{N}\!\left(x_t \,\middle|\, \alpha_t x_0,\; \sigma_t^2 I\right),
\end{equation}
其中 $\alpha_t$ 为缩放系数，$\sigma_t^2$ 为噪声方差，$I$ 为单位矩阵。
因此$x_t$可以由式\cref{eq:x_t_eq}表达
\begin{equation}\label{eq:x_t_eq}
x_t = \alpha_t x_0 + \sigma_t \epsilon,\qquad \epsilon \sim \mathcal{N}(0,I).
\end{equation}
在常见的扩散路径设定中，$\alpha_t$ 随 $t$ 单调递减而 $\sigma_t$ 随 $t$ 单调递增，从而使样本分布沿时间从 $p_{\mathrm{data}}$ 平滑过渡到先验噪声分布。
进一步地，在随机插值与流匹配一类设定中，常将时间范围限制为 $t \in [0,1]$，并施加端点条件
\[
\alpha_0 = 1,\ \sigma_0 = 0,\qquad \alpha_1 = 0,\ \sigma_1 = 1,
\]
于是 $x_{t=0}=x_0$ 且 $x_{t=1}=\epsilon$，从而 $x_t$ 在数据样本与标准高斯噪声之间形成精确的插值路径。
通常将 $t \in [0,1]$ 视为前向扩散过程的演化方向，即从数据逐步走向噪声先验，
相应地，生成过程对应其反向动力学，可写作 $t \in [1,0]$。
如\cref{eq:forward_diffusion}所示，系数 $\alpha_t$ 与 $\sigma_t$ 决定了时刻 $t$ 的扰动状态 $x_t$，并且二者均为时间 $t$ 的函数。因而，$\{\alpha_t,\sigma_t\}_{t\in[0,1]}$ 也可视为对前向扩散过程的一种路径参数化，常被称为扩散轨迹。在给定相同的起点样本 $x_0$ 与终点噪声 $\epsilon$ 的条件下，不同的系数函数将诱导不同的中间状态分布 $p_t$，从而对应不同的扩散轨迹。

本文采用广义的 VP-SDE 路径参数化，即选取三角函数系数来定义扩散路径，其形式为
\[
\alpha_t = \cos\!\left(\frac{\pi t}{2}\right), \qquad 
\sigma_t = \sin\!\left(\frac{\pi t}{2}\right).
\]
该选择满足恒等式 $\alpha_t^2+\sigma_t^2=1$，因此可使边缘方差随时间保持近似恒定，从而体现出方差保持的特性，可以在训练稳定性与生成质量之间保持较高的质量。

\xsubsection{流模型}{Flow Model}
与扩散模型通过随机扰动并学习其反向随机动力学不同，流模型通常以确定性连续时间动力系统为基础，通过学习一个时间依赖的速度场，将简单先验分布逐步变换为数据分布。设 $t\in[0,1]$ 表示变换时间，$x_t\in\mathbb{R}^d$ 为随时间演化的随机变量，流模型假设其轨迹满足如下常微分方程
\begin{equation}\label{eq:flow_ode}
\frac{\mathrm{d}x_t}{\mathrm{d}t} = v_\theta(x_t,t),
\end{equation}
其中 $v_\theta(\cdot,t)$ 为可学习的速度场，通常由神经网络参数化。由动力系统诱导的边缘分布 $p_t(x)$ 满足连续性方程
\begin{equation}\label{eq:continuity_eq}
\partial_t p_t(x) + \nabla\!\cdot\!\bigl(p_t(x)\,v_\theta(x,t)\bigr)=0,
\end{equation}
因此只要 $v_\theta$ 足够表达真实的传输场，便可在 $t=1$ 将先验分布 $p_1$ 精确推送到 $p_0\approx p_{\mathrm{data}}$。在生成时，通常从 $x_1\sim p_1$ 出发，沿时间反向求解 \cref{eq:flow_ode} 得到 $x_0$，从而实现样本生成。与离散化的归一化流模型相比，\cref{eq:flow_ode} 对应的连续归一化流（Continuous Normalizing Flow, CNF）提供了更自然的连续时间表述，其对数密度可通过伴随 ODE 或瞬时雅可比迹来更新
\begin{equation}\label{eq:cnf_logp}
\frac{\mathrm{d}}{\mathrm{d}t}\log p_t(x_t) = - \nabla\!\cdot v_\theta(x_t,t),
\end{equation}
从而在训练阶段可以最大化似然或最小化相关的 KL 目标。




